\bookchapter{פרק 19: אבחון תקלות והצלה עצמית}

\chapterintro{להגיד את האמת, אחרי שנה של שימוש הטלפון מתחיל להאט, זה נורמלי מאוד. הסיבה לא בהכרח שהטלפון התקלקל, לרוב מדובר ב“אבק” - קבצים זבל, תהליכי רקע, שטח אחסון שמתמלא לאט.}

\booksection{קודם אבחן, אל תמהר לתקן:}

בדוק שטח אחסון

\booksubsection{: לך להגדרות אחסון ובדוק שטח פנוי. אם נשאר פחות מ-10\%, המערכת תהיה איטית מאוד, כי היא נאבקת לפנות מקום. מניסיוני: שמור לפחות 15\%-20\% שטח פנוי, רק אז הטלפון זורם.}

\begin{enumerate}
  \item \textbf{בדוק דירוג צריכת חשמל}: בהגדרות סוללה בדוק את דירוג צריכת החשמל ב-24 השעות האחרונות. אם אפליקציה מסוימת תופסת נתח גדול אבל בקושי השתמשת בה, היא רצה ברקע. הסר או הגבל פעילות רקע.
  \item \textbf{חשוב על מה עשית לאחרונה}: האם האט אחרי עדכון מערכת? (אחרי עדכון גרסה גדולה האטה בימים הראשונים נורמלית, המערכת בונה מחדש אינדקס ברקע). או אחרי התקנת אפליקציה חדשה? (נסה להסיר).
  \item \textbf{הסתכל בסביבת השימוש}: טלפון בשמש, בנרתיק עבה, טעינה ומשחק במקביל - חימום בלתי נמנע. חום גורם ל-CPU להוריד תדר אוטומטית, והטלפון פתאום נראה איטי. תן לו להתקרר ונסה שוב.
  \item \textbf{שיטות טיפול (מהפשוט למורכב)}: הפעל מחדש את הטלפון (פותר 60\% מהבעיות) נקה שטח אחסון (מחק אפליקציות וסרטונים מיותרים) אם אפליקציה מסוימת מזיקה, נקה מטמון או הסר והתקן מחדש מוצא אחרון: גבה נתונים ואפס להגדרות יצרן.
\end{enumerate}

\textbf{שתיים, אפליקציה קורסת, לא מגיבה - האפליקציה לא ממושמעת}אפליקציה נסגרת בפתיחה, או קופאת באמצע שימוש. תהליך הטיפול הוא כמו ביקור רופא, צעד אחר צעד:

\booksection{קודם סגור ופתח מחדש}

: סגור את האפליקציה מהרקע, ופתח שוב. רוב הבאגים הזמניים נפתרים בזה.

\begin{enumerate}
  \item \textbf{עדכן את האפליקציה}: בדוק בחנות האפליקציות אם יש עדכון. לפעמים הגרסה ישנה מדי ולא תואמת למערכת.
  \item \textbf{בדוק הרשאות}: אפליקציות מסוימות דורשות הרשאה מסוימת כדי לעבוד, אבל לא נתת. למשל אפליקציית מצלמה בלי הרשאת מצלמה קורסת בפתיחה. בהגדרות מצא את האפליקציה, תן את כל ההרשאות הנדרשות.
  \item \textbf{נקה מטמון}: הגדרות ניהול אפליקציות מצא את האפליקציה אחסון נקה מטמון. שים לב: "נקה מטמון" לא "נקה נתונים"! ניקוי נתונים ימחק את מצב ההתחברות וההגדרות, כמו איפוס יצרן לאפליקציה.
  \item \textbf{הסר והתקן מחדש}: המוצא האחרון. לפני הסרה אשר שהנתונים לא יאבדו - ב-WeChat גבה היסטוריית צ'אט, במשחקים אשר שהחשבון קשור למספר טלפון.
  \item \textbf{שלוש, לא טוען או טוען לאט - אולי לא הסוללה אשמה}"לא טוען" היא בעיה שגורמת לחרדה, אבל לרוב הסיבה היא באביזרים, לא בטלפון עצמו.
\end{enumerate}

\booksection{קודם בדוק את פתח הטעינה}

: בעזרת קיסם או כלי קטן לא מוליך, נקה בעדינות אבק ופתיתים מפתח הטעינה. זה קרה לי - חוט קטן מהכיס נתקע בדיוק בפתח, חשבתי שאין מגע, הוצאתי והכל עבד.

\begin{itemize}
  \item \textbf{נסה כבל ומטען אחרים}: כבלים ניזוקים יותר ממטענים, במיוחד כבלים שמתקפלים הרבה. החלף כבל ומטען ידועים כתקינים ובדוק, תוכל לקבוע מהר מה התקלקל.
  \item \textbf{בעיות טעינה אלחוטית}: אם משתמש בטעינה אלחוטית, בדוק שהנרתיק לא עבה מדי, או שאין פיסת מתכת מאחוריו (כמו תושבת מגנטית). גם, הנח את הטלפון בדיוק במרכז משטח הטעינה - אם מוטה אולי לא יטען.
  \item \textbf{טעינה איטית}: ודא שהספק המטען מספיק (iPhone לטעינה מהירה דורש לפחות מטען PD 20W), שהכבל תומך בטעינה מהירה (לא כל כבל מעביר זרם 3A+), שהטלפון לא חם מדי (חום יוריד הספק להגנת הסוללה).
  \item \textbf{ארבע, לא מתחבר לרשת, Wi-Fi מתנתק כל הזמן}בחוץ הסלולר עובד טוב, בבית לא מתחבר ל-Wi-Fi; או מחובר אבל אין אינטרנט. בדוק לפי הסדר:
\end{itemize}

\booksection{מתג מצב טיסה}

: משוך את מרכז הבקרה, הפעל מצב טיסה ל-5 שניות וכבה. זה גורם לטלפון לחפש רשת מחדש, יעיל מאוד.

\begin{enumerate}
  \item \textbf{"שכח" Wi-Fi וחבר מחדש}: בהגדרות Wi-Fi, לחץ על סימן הקריאה ליד שם הרשת (באנדרואיד גלגל שיניים), בחר "התעלם מרשת זו" או "שכח", ואז הזן סיסמה מחדש.
  \item \textbf{הפעל מחדש את הנתב}: הרבה נתבים בבית לא כובים חודשים, עם הזמן נעשים לא יציבים. נתק מהחשמל ל-30 שניות וחבר שוב, רוב בעיות ה-Wi-Fi נפתרות.
  \item \textbf{בדוק אם הטלפון אשם}: האם מכשירים אחרים מתחברים לאותו Wi-Fi בסדר? אם כן ורק הטלפון שלך לא, אולי צריך להפעיל מחדש או לאפס הגדרות רשת (הגדרות כללי העברה או איפוס איפוס הגדרות רשת, לא מוחק נתונים, אבל מוחק את כל סיסמאות ה-Wi-Fi).
  \item \textbf{חמש, בעיות מסך: מגע לא תקין, הבהוב, שאריות תמונה}מגע לא רגיש
\end{enumerate}

\booksection{: קודם הסר את המגן המחוזק ובדוק אם המגן אשם - מגנים זולים אחרי זמן באמת פוגעים במגע. מים על המסך או שמן על היד גם גורמים לסטיית מגע. אם לא, הפעל מחדש ונסה שוב.}

\begin{itemize}
  \item \textbf{שאריות/תמונת רפאים על המסך (burn-in)}: הצגת אותה תמונה זמן רב (למשל השארת המסך דלוק תוך גלילה ב-Douyin), תשאיר "תמונת רפאים" על מסך OLED. זה לא שהמסך התקלקל, רק זמני. כיבוי המסך לחצי שעה בדרך כלל משחזר.
  \item \textbf{מסך מהבהב}: אם ההבהוב הוא בבהירות נמוכה, זה עשוי להיות בעיית PWM dimming, נפוץ במסכי OLED. בהגדרות חפש "DC dimming" או "anti-flicker", הפעלה תעזור הרבה. לא לכל הטלפונים יש את הפונקציה הזו.
  \item \textbf{אזור מסוים במסך שחור/מואר לגמרי}: זה חמור יותר, כנראה נזק פיזי למסך, צריך החלפת מסך. מסך שהוחלף בתיקון לא רשמי, הצבע והבהירות בדרך כלל נחותים מהמקורי, וגם פוגע בעמידות למים.
  \item \textbf{שש, לא נדלק, נתקע בלוגו - אל תבהל}מצב כזה גורם לדפיקות לב מהירות, אבל תירגע, קודם נסה את אלה:
\end{itemize}

\booksection{תטען חצי שעה ואז נסה}

: לרוב זה פשוט שהסוללה התרוקנה לגמרי. חבר מטען מקורי, חכה לפחות 30 דקות, ואז נסה להדליק. חלק מהטלפונים אחרי ריקון מוחלט צריכים להגיע לרמה מסוימת כדי להגיב.

\begin{itemize}
  \item \textbf{הפעלה מאולצת}: iPhone full-screen: לחץ פעם אחת על הגדלת עוצמה, פעם אחת על הנמכת עוצמה, לחץ ארוך על כפתור הצד עד שמופיע לוגו Apple (אולי צריך להחזיק 10-15 שניות). אנדרואיד שיטה כללית: לחץ ארוך על כפתור ההפעלה מעל 10 שניות, עד לרטט או הדלקת מסך.
  \item \textbf{ניקוי פתח טעינה}: מצב כזה נקרא "Apple לבן" (iPhone) או "נתקע בלוגו" (אנדרואיד). אם הפעלה מאולצת לא עוזרת, אולי צריך להיכנס למצב שחזור ולתקן עם מחשב. iPhone חבר למחשב פתח iTunes/Finder, אחרי כניסה למצב שחזור בחר "עדכן" (לא "שחזר"), יתקין מערכת מחדש אך ישמור נתונים. אנדרואיד כנס ל-Fastboot או Recovery, כל מותג שילוב מקשים שונה, מומלץ לחפש דגם הטלפון שלך + "כניסה למצב שחזור".
  \item \textbf{אם כל הנ"ל לא עוזר}: אז רק לפנות לתיקון מקצועי. אם גובו נתונים, לא צריך לדאוג יותר מדי.
  \item \textbf{שבע, בעיות קול: רמקול לא עובד, קול שיחה חלש}רמקול ללא קול
\end{itemize}

\booksection{: קודם ודא שלא מצב שקט (מתג השקט בצד iPhone), מצב לא להפריע, או שהעוצמה במינימום. נשמע טיפשי, אבל זה 40\% מהמקרים. אז בדוק אם Bluetooth מחובר לאוזניות או רמקול.}

\begin{itemize}
  \item \textbf{בשיחה הצד השני לא שומע אותך}: בדוק שפתח המיקרופון לא סתום - החורים הקטנים בתחתית הטלפון בקלות נסתמים באבק. נקה בעדינות במברשת רכה. אל תשתמש במחט או קיסם, תקלקל את רשת האבק.
  \item \textbf{קול צרוד/מתפצל}: אם הרמקול נרטב - אל תשתמש במייבש שיער חם! ספוג נוזלים במגבת נייר, ואז הנח לייבוש טבעי. באינטרנט אומרים שאפליקציות הוצאת מים בסאונד עוזרות (משמיעות צליל תדר נמוך להוציא מים), אפשר לנסות.
  \item \textbf{שמונה, תיקון רשמי, הימנע מהמלכודות}אם ניסית את כל השיטות למעלה ולא עזר, צריך תיקון מקצועי. אבל שוק התיקונים עמוק, אספר על המלכודות שנתקלתי בהן והניסיון שלי:
\end{itemize}

\booksection{מצא רשמי אמיתי, לא “נראה כמו רשמי”}

: באתר המותג (Apple.com.cn, Huawei.com.cn וכו') חפש רשימת “ספקים מורשים”. חנות ברחוב עם לוגו Apple לא מורשית, עשויה להשתמש בחלקים לא מקוריים, אחרי התיקון עמידות המים נעלמת.

\begin{itemize}
  \item \textbf{לפני שליחה לתיקון חובה לעשות}: גבה את כל הנתונים (אל תסמוך שלא יאבדו בתיקון), צא מכל החשבונות החשובים (Apple ID, WeChat, Alipay), כבה "איתור הטלפון שלי". חלק מהמתקנים יבקשו את קוד פתיחת המסך, זה סיכון פרטיות - עדיף לצאת מהחשבונות קודם בעצמך.
  \item \textbf{מלכודת החלפת מסך}: מסכים לא רשמיים מחולקים ל"מסך מקורי לחוץ" (זכוכית מקורית מוחלפת), "מסך לחוץ אחורה" (זכוכית לא מקורית), "מסך מורכב" (הכל לא מקורי). המחירים שונים פי כמה, אבל גם איכות התצוגה והעמידות שונות מאוד. אם הטלפון באחריות, לך ברשמי.
  \item \textbf{מלכודת החלפת סוללה}: סוללות לא מקוריות עשויות להגזים בקיבולת, ואחרי חצי שנה דועכות יותר מסוללה מקורית אחרי שנתיים. מחיר החלפת סוללה ב-Apple רשמי לא יקר הרבה יותר מלא רשמי, מומלץ ברשמי.
  \item \textbf{מלכודת שחזור נתונים}: אחרי נזק מים/נפילה, אם בתיקון אומרים "שחזור נתונים בתשלום נוסף", זהירות - חלק מהעסקים הלא הגונים יגזימו בנזק כדי לגבות יותר. שחזור נתונים רשמי הוא קודם הצעת מחיר ואז ביצוע, ואם לא מצליח לא משלמים.
  \item \textbf{הפתרון האוניברסלי הראשון: הפעלת הטלפון מחדש}לא משנה איזו בעיה נתקלת - קפיאה, קריסה, לא מתחבר לרשת, קול חריג - קודם הפעל את הטלפון מחדש. זה נשמע כמו זלזול, אבל מניסיוני, לפחות 60\% מהתקלות הקטנות נפתרות אחרי הפעלה מחדש. הטלפון, כמו מחשב, צובל לאורך זמן מיני דברים מבולגנים בזיכרון, הפעלה מחדש היא איפוס. לכן, בפעם הבאה שהטלפון מתקלקל, אל תמהר לחפש "מה לעשות", קודם הפעל מחדש ובדוק.
\end{itemize}

\begin{tcbwarn}{זכור עכשיו: פעולות הצלה אלה צריכות להיות מוכרות לך היטב}
הפעלה מאולצת

\end{tcbwarn}

\begin{tcbpractice}{: iPhone full-screen: הגדל עוצמה הנמך עוצמה לחץ ארוך על כפתור ההפעלה עד ללוגו. אנדרואיד: לחץ ארוך על כפתור ההפעלה 10-15 שניות עד רטט.}
\begin{enumerate}
  \item \textbf{ניקוי פתח טעינה}: מצא כלי קטן (קיסם או מחט SIM), נקה בעדינות אבק מפתח הטעינה.
  \item \textbf{חיבור Wi-Fi מחדש}: בהגדרות Wi-Fi "שכח" את הרשת הנוכחית, הזן סיסמה מחדש.
  \item \textbf{מתג מצב טיסה}: משוך מרכז בקרה, הפעל מצב טיסה, ספור 5 שניות, כבה, בדוק אם הרשת חזרה.
  \item \textbf{צלם את השלבים של הפעולות האלה, או כתוב בפתק. בזמן תקלה לא צריך להילחץ ולחפש, פשוט פתח את הפתק.}פרק 20: מצב מפתח והרשאות מתקדמות
  \item פרק זה מיועד למשתמשים שרוצים לחקור לעומק את "הפונקציות הנסתרות" של הטלפון. מצב מפתח, הרשאות Root, נעילת BL... מושגים שנשמעים גיקיים, אבל למעשה לא רחוקים מהמשתמש הרגיל.
\end{enumerate}

\end{tcbpractice}
